\section{Architektura frameworku}\label{sec:architektura}

Ten rozdział pokazuje całość: jakie kroki składają się na weryfikację migracji, jakie
pliki przechodzą między krokami, które komponenty je produkują, w których repozytoriach
mieszkają i jakie zasady projektowe obowiązują wszystkie z nich. Szczegóły każdego
komponentu są w rozdziałach~\ref{sec:rdzen}--\ref{sec:bench}.

\subsection{Łańcuch weryfikacji}\label{sec:architektura-lancuch}

Łańcuch ma siedem kroków. Każdy kończy się plikiem, który czyta następny krok.

\begin{tabularx}{\linewidth}{@{}p{2.3cm} >{\raggedright\arraybackslash}X >{\raggedright\arraybackslash}p{3.3cm}@{}}
\toprule
Krok & Komponent & Wynik \\
\midrule
1 Capture & C1 --- proxy nagrywające przed systemem legacy, bez zmiany kodu & \repopath{*.skcap} \\
2 Contract & C2 (miner, faza 02) i człowiek; C3 --- silnik kontraktu & \repopath{contract.yaml} \\
3 Standards & C4 --- standardy dostarczane agentowi migracyjnemu jako skills & pull request agenta \\
4 Gate & C5 --- analyzery Roslyn na zmienionych liniach (Portcullis) & \repopath{*.sarif} \\
5 Replay & C6 --- obie wersje w identycznych warunkach; C7 porównuje pole po polu względem kontraktu & \repopath{*.skrun}, \repopath{verdict.json} \\
6 Mutations & C9 --- wstrzyknięte defekty, które kontrakt musi wykryć & kill rate per klauzula \\
7 Evidence & C10 --- jeden pakiet na pull request & \repopath{evidence/} \\
\bottomrule
\end{tabularx}

\begin{figure}[htbp]
  \centering
  \includegraphics[width=\linewidth,height=0.62\textheight,keepaspectratio]{\dgm{a2-lancuch-weryfikacji}}
  \caption{Łańcuch weryfikacji i artefakty między krokami.}
  \label{fig:lancuch}
\end{figure}

Strona legacy i kandydat są uruchamiane, ale nie są modyfikowane przez łańcuch: capture
stoi przed systemem legacy jako proxy, a replay wysyła nagrane żądania do obu wersji.
Wszystko, co następuje po replay, działa na plikach.

\zrodla{\path{letsgolegacy.secondkey/README.md},
\path{letsgolegacy.secondkey/docs/WORKPLAN.md}}

\subsection{Artefakty są interfejsem}\label{sec:architektura-artefakty}

\finding{Każdy komponent czyta i pisze pliki, nigdy wnętrza innego komponentu.} Dowolny
komponent można wymienić, dopóki czyta i pisze te same artefakty. Dlatego formaty są
wersjonowane, opisane schematami JSON i pilnowane testami na próbkach
(rozdział~\ref{sec:formaty}).

\begin{xltabular}{\linewidth}{@{}>{\raggedright\arraybackslash}p{2.9cm} >{\raggedright\arraybackslash}X >{\raggedright\arraybackslash}X >{\raggedright\arraybackslash}X@{}}
\toprule
Artefakt & Format & Produkuje & Czyta \\
\midrule
\endfirsthead
\toprule
Artefakt & Format & Produkuje & Czyta \\
\midrule
\endhead
\repopath{*.skcap} & JSONL, jedno zdarzenie na linię, korelowane identyfikatorem wymiany (\emph{exchange id}) & C1 capture & C2 miner, C6 replay \\
\repopath{contract.yaml} & YAML walidowany JSON Schema & C2 miner i człowiek & C3 silnik, C7 komparator, C9 mutacje \\
\repopath{standards/} & Markdown → skills i NuGet & C4 (\path{letsgolegacy.secondkey-standards}) & agent migracyjny, C5 bramka \\
\repopath{*.sarif} & SARIF 2.1 & C5 bramka (\path{letsgolegacy.portcullis}) & C10 evidence \\
\repopath{*.skrun} & JSONL, wyniki obu stron dla każdego scenariusza & C6 replay & C7 komparator \\
\repopath{verdict.json} & JSON & C7 komparator, C9 mutacje & C10 evidence, C11 portal \\
\repopath{evidence.dsse} & in-toto / DSSE & C10 evidence (podpis od fazy 03) & C11 portal, audytor \\
\bottomrule
\end{xltabular}

Konsekwencja widoczna w benchu: granica między maszynami może przebiegać wszędzie tam,
gdzie da się skopiować plik. Tylko capture i replay muszą mieć dostęp HTTP do działającego
systemu; kontrakt, porównanie i pakiet dowodowy to „plik na wejściu, plik na wyjściu”
(sekcja~\ref{sec:architektura-topologia}).

\zrodla{\path{letsgolegacy.secondkey/docs/WORKPLAN.md},
\path{letsgolegacy.secondkey/docs/adr/0001-artifacts-are-the-interface.md},
\path{letsgolegacy.secondkey/schemas/skcap.schema.json},
\path{letsgolegacy.secondkey-nopcommerce-bench/docs/TOPOLOGY.md}}

\subsection{Komponenty C0--C12 i repozytoria}\label{sec:architektura-komponenty}

\begin{xltabular}{\linewidth}{@{}l >{\raggedright\arraybackslash}X >{\raggedright\arraybackslash}p{5.2cm} l@{}}
\toprule
Id & Komponent & Repozytorium & Faza \\
\midrule
\endfirsthead
\toprule
Id & Komponent & Repozytorium & Faza \\
\midrule
\endhead
C0 & CLI \repopath{sk} & \path{letsgolegacy.secondkey} & 01 \\
C1 & capture HTTP (proxy YARP, tylko ruch przychodzący) & \path{letsgolegacy.secondkey} & 01 \\
C1 & capture SQL, kolejek, tokenizacja PII & \path{letsgolegacy.secondkey-capture} & 02 \\
C2 & miner kontraktu (\repopath{sk mine}) & \path{letsgolegacy.secondkey} & 02 \\
C3 & silnik kontraktu & \path{letsgolegacy.secondkey} & 01 \\
C4 & pakiet standardów & \path{letsgolegacy.secondkey-standards} & 01 \\
C5 & bramka na PR (Portcullis) & \path{letsgolegacy.portcullis} & 01 \\
C6 & replay i reset stanu & \path{letsgolegacy.secondkey} & 01 \\
C7 & komparator i werdykt & \path{letsgolegacy.secondkey} & 01 \\
C8 & doradcza triage różnic (nigdy nie zmienia werdyktu) & \path{letsgolegacy.secondkey} & 02--03 \\
C9 & silnik mutacji (\repopath{sk mutate}) & \path{letsgolegacy.secondkey} & 02 \\
C10 & pakiet dowodowy (bez podpisu w fazie 01, podpis w fazie 03) & \path{letsgolegacy.secondkey} & 01, 03 \\
C11 & portal dowodów & \path{letsgolegacy.secondkey-portal} & 03 \\
C12 & pakowanie w granicy klienta (sandbox) & \path{letsgolegacy.secondkey-sandbox} & 02 \\
\bottomrule
\end{xltabular}

\begin{figure}[htbp]
  \centering
  \includegraphics[width=\linewidth,height=0.7\textheight,keepaspectratio]{\dgm{a1-mapa-repozytoriow}}
  \caption{Repozytoria frameworku letsgolegacy i to, co przepływa między nimi. Nazwy bez wspólnego prefiksu \texttt{letsgolegacy.}}
  \label{fig:mapa}
\end{figure}

W fazie 01 \repopath{sk mine} i \repopath{sk mutate} istnieją jako stuby, które kończą
się kodem 70 (\repopath{NotImplemented}); kontrakt pisze człowiek.

\zrodla{\path{letsgolegacy.secondkey/README.md},
\path{letsgolegacy.secondkey/docs/WORKPLAN.md},
\path{letsgolegacy.secondkey/docs/cli.md},
\path{letsgolegacy.secondkey-capture/docs/WORKPLAN.md},
\path{letsgolegacy.secondkey-sandbox/docs/WORKPLAN.md},
\path{letsgolegacy.secondkey-portal/docs/WORKPLAN.md}}

\subsection{Zasady projektowe}\label{sec:architektura-zasady}

Rdzeń egzekwuje siedem zasad projektowych. Pozostałe repozytoria się na nie powołują.

\begin{enumerate}
  \item \textbf{Werdykt jest deterministyczny.} O pass/fail decyduje wyłącznie kod bez
        modelu językowego. Model może szkicować klauzule kontraktu albo porządkować
        różnice do przeglądu, ale nigdy nie decyduje. Nic w \repopath{SecondKey.Contract}
        ani \repopath{SecondKey.Compare} nie woła modelu językowego; ten sam kontrakt i
        ten sam przebieg dają zawsze ten sam werdykt (poza polem \repopath{createdAt}).
  \item \textbf{Wszystko działa w granicy klienta; domyślnie zero ruchu wychodzącego.}
        Repozytorium \path{letsgolegacy.secondkey-sandbox} pakuje łańcuch dla różnych
        rodzajów takiej granicy.
  \item \textbf{Tylko technologie, które odbiorca już ma}: .NET, SQL Server, Entra ID,
        GitHub albo Azure DevOps, Kubernetes.
  \item \textbf{System legacy jest nietykalny}: capture nie wymaga zmiany kodu ani
        rekompilacji. Czas i losowość w kodzie legacy (\repopath{DateTime.Now},
        \repopath{Guid.NewGuid()}) obsługuje normalizacja w komparatorze, a nie zmiana
        systemu legacy.
  \item \textbf{Dowód jest artefaktem kryptograficznym, nie dashboardem}: pakiet zawiera
        skróty SHA-256 wszystkich plików i niepodpisany statement in-toto; podpis dochodzi
        w fazie 03 bez zmiany formatu.
  \item \textbf{Brak własnej usługi tożsamości}; używany jest IdP klienta.
  \item \textbf{Każdy krok działa z CLI w pipeline}; portal jest oknem na dowody, a nie
        interfejsem.
\end{enumerate}

Obok siedmiu zasad cały łańcuch wiążą trzy decyzje:

\begin{itemize}
  \item \textbf{Różnica bez wyjaśnienia zawodzi (fail closed).} Różnica, której nie tłumaczy
        żadna klauzula ani reguła normalizacji, jest regresją \repopath{uncovered-difference}
        i blokuje, dopóki przeglądnięta zmiana kontraktu nie doda reguły normalizacji
        (różnica nie ma znaczenia) albo klauzuli (ma znaczenie). Strona bez odpowiedzi to
        regresja \repopath{missing-result} (ADR 0003, decyzja D-2 analizy fazy 01).
  \item \textbf{Kontrakt nie czyta prozy i nie przechodzi pusto.} Wartości ze stron HTML
        wyciągają nazwane ekstraktory z jawnym typem i kulturą; klauzula, do której nie
        trafiło żadne żądanie, jest \repopath{unexercised}, nigdy spełniona (ADR 0002).
  \item \textbf{Agent i bramka dostają te same standardy.} Standardy są spisane raz
        i kompilowane do skilla agenta oraz do konfiguracji analyzerów bramki; rozbieżność
        między tym, co usłyszał agent, a tym, co egzekwuje bramka, jest błędem CI, a nie
        znaleziskiem (rozdział~\ref{sec:standardy}).
\end{itemize}

Decyzje ADR rdzenia opisuje podrozdział~\ref{sec:rdzen-adr}.

\zrodla{\path{letsgolegacy.secondkey/docs/WORKPLAN.md},
\path{letsgolegacy.secondkey/CLAUDE.md},
\path{letsgolegacy.secondkey/docs/formats/verdict.md},
\path{letsgolegacy.secondkey/docs/analysis/phase-01.md},
\path{letsgolegacy.secondkey/docs/adr/0002-contract-language.md},
\path{letsgolegacy.secondkey/docs/adr/0003-verdict-classes.md},
\path{letsgolegacy.secondkey-standards/docs/WORKPLAN.md},
\path{letsgolegacy.secondkey-portal/docs/WORKPLAN.md}}

\subsection{Fazy i bramka fazy 01}\label{sec:architektura-fazy}

Praca jest podzielona na fazy. Faza 01 to demonstracja; jej bramką jest publiczny pakiet
dowodowy na nopCommerce 3.x, czyli ticket P9 benchu. Ścieżka krytyczna fazy 01:
S0 → S1 → S3--S6 → S7--S14 → S15 → (bench P3) → (bench P4--P8) → (bench P9). Faza 02
zaczyna się dopiero po bramce fazy 01.

\begin{xltabular}{\linewidth}{@{}l >{\raggedright\arraybackslash}X@{}}
\toprule
Faza & Zakres według planów pracy repozytoriów \\
\midrule
\endfirsthead
\toprule
Faza & Zakres według planów pracy repozytoriów \\
\midrule
\endhead
01 & Rdzeń: CLI, capture HTTP, silnik kontraktu, replay z resetem snapshotem SQL Server,
komparator, niepodpisany pakiet dowodowy, job end-to-end (S0--S15; S16 przeniesiony do
fazy 02). Portcullis: import, zmiana nazwy, reguły migracyjne, SARIF z filtrem zmienionych
linii i baseline (R0--R3); wsparcie P8 (reguły Portcullis na pull requeście kandydata
benchu) zaplanowane. Standardy: format, generator, drift check, skill dla modernize-dotnet
(C4-a--c, R4). Bench: P1--P10. legacy-risk-scan: L18. Sandbox: decyzja o topologii
demonstracyjnej D7 w stanie „proposed”, zrobiona, gdy ADR zostanie przyjęty. \\
02 & \repopath{sk mine} (miner kontraktu), \repopath{sk mutate} (silnik mutacji), capture
SQL, kolejek i PII w \path{secondkey-capture}, triage doradcza C8, pakowanie sandbox (Helm,
Hyper-V, Bicep, bundle air-gapped), reguły warstw ArchUnitNET w Portcullis, standardy
klienta jako skills, kanały wydań standardów przez NuGet i tagi git z wersją przypiętą
w każdym repozytorium (02--03), outbound stubs (S16). \\
03 & Podpis pakietu (in-toto w kopertach DSSE, cosign z kluczem klienta), biblioteka
fragmentów kontraktów, portal dowodów C11, paczki reguł per estate, pakiet Azure Managed
Application. \\
04 & Java: capture usług Java, analyzery Java w Portcullis (modernize-java). \\
\bottomrule
\end{xltabular}

\zrodla{\path{letsgolegacy.secondkey/docs/WORKPLAN.md},
\path{letsgolegacy.portcullis/docs/WORKPLAN.md},
\path{letsgolegacy.secondkey-standards/docs/WORKPLAN.md},
\path{letsgolegacy.secondkey-nopcommerce-bench/docs/WORKPLAN.md},
\path{letsgolegacy.secondkey-capture/docs/WORKPLAN.md},
\path{letsgolegacy.secondkey-sandbox/docs/WORKPLAN.md},
\path{letsgolegacy.secondkey-portal/docs/WORKPLAN.md}}

\subsection{Topologia: Windows dla legacy, Linux dla reszty}\label{sec:architektura-topologia}

System legacy na .NET Framework potrzebuje Windows: nopCommerce 3.90 to ASP.NET MVC 5
serwowany przez IIS z bazą SQL Server. Łańcuch Second Key jest napisany na .NET 10 i
działa wszędzie. Demonstracyjna topologia (decyzja D7, w stanie „proposed”) umieszcza więc
stronę legacy na jednym hoście Windows, a resztę łańcucha na Linuksie. Tylko dwa kroki
muszą sięgać do działającego sklepu przez HTTP: capture i replay. Działają obok IIS na
hoście Windows, a ich wyniki (\repopath{*.skcap}, \repopath{*.skrun}) wychodzą jako
artefakty CI do jobu linuksowego, który liczy werdykt i pakiet. Hostowane runnery GitHub
nie współdzielą sieci, więc job linuksowy i tak nie mógłby sięgnąć do IIS na runnerze
Windows.

\begin{figure}[htbp]
  \centering
  \includegraphics[width=\linewidth,height=0.6\textheight,keepaspectratio]{\dgm{a7-topologia-bench}}
  \caption{Topologia benchu: host Windows ze stroną legacy, capture i replay; Linux
           dla werdyktu i pakietu.}
  \label{fig:topologia}
\end{figure}

\zrodla{\path{letsgolegacy.secondkey-nopcommerce-bench/docs/TOPOLOGY.md},
\path{letsgolegacy.secondkey-sandbox/docs/WORKPLAN.md}}

\subsection{Znane trudne problemy}\label{sec:architektura-problemy}

Plan pracy rdzenia wymienia problemy, które są śledzone, a nie ukrywane:

\begin{itemize}
  \item \textbf{C1:} przypisanie instrukcji SQL do żądania HTTP bez zmian w kodzie to
        heurystyka z jawnie raportowanym marginesem niepewności.
  \item \textbf{C2:} fałszywe niezmienniki z małych prób; miner raportuje wsparcie
        statystyczne dla każdej klauzuli i milczy przy zbyt małej liczbie danych.
  \item \textbf{C6:} czas i losowość w kodzie legacy bez wstrzykiwania zależności są
        obsługiwane normalizacją w C7, a nie ingerencją w system legacy.
  \item \textbf{C9/C5:} czas działania na dużych solution; mutacje i analiza tylko kodu
        zmienionego w PR.
  \item \textbf{C10:} postać raportu, którą zaakceptuje audytor, jest ustalana z
        pierwszymi partnerami, a nie zgadywana.
\end{itemize}

\zrodla{\path{letsgolegacy.secondkey/docs/WORKPLAN.md}}
